\subsection{逼近空间的例子}

回顾，每个希尔伯特空间都是逼近空间（III, p. 15, 命题 10）。本节给出其它一些例子。

若 X 是局部紧空间，我们用 \(\mathcal{C}_{0}(\mathrm{X};\mathrm{K})\)，或简记作 \(\mathcal{C}_{0}(\mathrm{X})\)，表示在无穷远处消失的从 X 到 K 的连续函数所构成的巴拿赫空间，赋以由 \(\|f\|=\sup_{x\in\mathrm{X}}|f(x)|\)（对 \(f\in\mathcal{C}_{0}(\mathrm{X})\)）定义的范数。当 X 紧时，这个空间与从 X 到 K 的连续函数构成的空间 \(\mathcal{C}(\mathrm{X})\) 重合。

引理 5. --- 设 X 是紧拓扑空间，F 是 \(\mathcal{C}(\mathrm{X})\) 的一个有限子集，\(\varepsilon > 0\) 是实数。

存在一个有限子集 \(X_{0} \subset X\) 和一个线性映射 \(u: \mathcal{C}(X_{0}) \to \mathcal{C}(X)\)（范数 \(\leqslant 1\)），使得 \(\|u(f|X_{0}) - f\| \leqslant \varepsilon\) 对一切 \(f \in F\) 成立。

由于集合 F 是 \(\mathcal{C}(\mathrm{X})\) 的等度连续子集，存在 X 的一个有限覆盖 \((\mathrm{U}_{i})_{i\in\mathrm{I}}\)，使得对每个 \(i\in I\) 和每个 \(f\in F\)，\(f(\mathrm{U}_{i})\) 的直径 \(\leqslant\varepsilon\)。对每个 \(i\in I\)，选取一个点 \(x_{i}\)（属于 \(U_{i}\)），并用 \(X_{0}\) 表示诸 \(x_{i}\)（\(i\in I\)）所成的有限集。

设 \((\varphi_{i})_{i\in\mathrm{I}}\) 是从属于 \((\mathrm{U}_{i})_{i\in\mathrm{I}}\) 的连续单位分解。如下定义一个从 \(\mathcal{C}(\mathrm{X}_{0})\) 到 \(\mathcal{C}(\mathrm{X})\) 中的线性映射 u：令

\[
u (g) = \sum_ {i \in \mathrm{I}} g (x _ {i}) \varphi_ {i}
\]

其中 \(g \in \mathcal{C}(X_{0})\)。线性映射 u 的范数 \(\leqslant 1\)，且对 F 中的一切 f 满足 \(\|u(f|X_{0}) - f\| \leqslant \varepsilon\)，由此即得断言。

命题 13. --- 设 X 是局部紧拓扑空间。巴拿赫空间 \(\mathcal{C}_{0}(\mathrm{X})\) 是逼近空间。

首先假设 X 紧。设 \(A \subset \mathcal{L}^{\mathrm{f}}(\mathcal{C}(\mathrm{X}))\) 是形如 \(f \mapsto u(f|X_{0})\) 的映射所成的集合，其中 \(X_{0}\) 是 X 的有限子集，而 \(u : \mathcal{C}(X_{0}) \to \mathcal{C}(X)\) 是范数 \(\leqslant 1\) 的线性映射。集合 A 等度连续。由引理 5，\(\mathcal{C}(X)\) 的恒等映射对 \(\mathcal{C}(X)\) 上的逐点收敛拓扑附着于 A。于是命题由 III, p. 14 的注记 3 得出。

转到一般情形。设 Y 是 X 的亚历山德罗夫紧化，\(\omega\) 是 Y 的无穷远点（TG, I, p. 67）。把 \(\mathcal{C}_{0}(X)\) 与 \(\mathcal{C}(Y)\) 中在 \(\omega\) 处消失的元素所成的集合等同起来。那么 \(\mathcal{C}(Y)\) 是 \(\mathcal{C}_{0}(X)\) 与 Y 上常值函数的向量空间的拓扑直和。由于按前面所证 \(\mathcal{C}(Y)\) 是逼近空间，故空间 \(\mathcal{C}_{0}(X)\) 是逼近空间（III, p. 14 的注记 1）。

推论. --- 每个交换星代数都是逼近空间。

事实上，这样的代数同构于一个局部紧空间上在无穷远处消失的连续函数所构成的代数（I, p. 108, 定理 1）。

注记. --- 若 E 是无穷维希尔伯特空间，则星代数 \(\mathcal{L}(\mathrm{E})\) 不具有逼近性质（A. SZANKOWSKI, \(\mathcal{B}(\mathrm{H})\) does not have the approximation property, Acta Math. 147 (1981), p. 89--108）。

引理 6. --- 设 X 是局部紧拓扑空间。设 \(\mu\) 是 X 上的一个正测度，\(p \in [1, +\infty[\)。设 F 是 \(\mathrm{L}_{\mathrm{K}}^{p}(\mathrm{X}, \mu)\) 的一个有限子集，\(\varepsilon > 0\) 是实数。

存在 \(\mathrm{L}_{\mathrm{K}}^{p}(\mathrm{X},\mu)\) 的一个范数 \(\leqslant 1\) 的有限秩投影子 u，使得对 F 中的每个 f，\(\|u(f)-f\|\leqslant\varepsilon\)。

设 \(\mathcal{P}\) 是 X 的有限分割 \(\pi = (\mathrm{K}_1, \ldots, \mathrm{K}_n, \mathrm{H})\) 所成的集合，其中 \(n \geqslant 1\) 是整数，而 \(\mathrm{K}_1, \ldots, \mathrm{K}_n\) 是 X 的测度非零的可积子集。对每个分割 \(\pi = (\mathrm{K}_1, \ldots, \mathrm{K}_n, \mathrm{H}) \in \mathcal{P}\)，如下定义自同态 \(v_\pi\)（\(\mathcal{L}^p(\mathrm{X}, \mu)\) 的）：令

\[
v _ {\pi} (f) = \sum_ {i = 1} ^ {n} \frac {1}{\mu (\mathrm{K} _ {i})} \left(\int_ {\mathrm{K} _ {i}} f d \mu\right) \varphi_ {\mathrm{K} _ {i}},
\]

其中 \(f \in \mathcal{L}_{\mathrm{K}}^{p}(\mathrm{X}, \mu)\)，\(\varphi_{\mathrm{K}_i}\) 是 \(\mathrm{K}_i\) 的特征函数。映射 \(v_\pi\) 通过过渡到商，诱导出一个投影子 \(u_\pi\)（\(\mathrm{L}_{\mathrm{K}}^{p}(\mathrm{X}, \mu)\) 的）。容易验证，\(u_\pi\) 的像是由满足以下条件的函数 \(f \in \mathcal{L}_{\mathrm{K}}^{p}(\mathrm{X}, \mu)\) 的类所构成的空间：\(f\) 在 H 上为零，且在 \(\mathrm{K}_i\) 上取常数（对 \(1 \leqslant i \leqslant n\)）。

我们来证明 \(\| u_{\pi}\| \leqslant 1\)。对 \(f\in \mathcal{L}_{\mathrm{K}}^{p}(\mathrm{X},\mu)\)，我们有

\[
\| u _ {\pi} (f) \| _ {p} ^ {p} = \sum_ {i = 1} ^ {n} \mu (\mathrm{K} _ {i}) ^ {1 - p} \left| \int_ {\mathrm{K} _ {i}} f d \mu \right| ^ {p}.
\]

由赫尔德不等式（INT, IV, p. 208, § 6, 第 4 目, 定理 2），对每个 i 我们有不等式

\[
\left| \int_ {\mathrm{K} _ {i}} f d \mu \right| ^ {p} \leqslant \mu (\mathrm{K} _ {i}) ^ {p - 1} \int_ {\mathrm{K} _ {i}} | f | ^ {p} d \mu ,
\]

由此得 \(\| u_{\pi}(f)\| _p\leqslant \| f\| _p.\)

设 \(\mathcal{E}\) 是 X 上只取有限多个值的可积函数在 \(\mathrm{L}_{\mathrm{K}}^{p}(\mathrm{X},\mu)\) 中的类所成的集合。由于 \(\mathcal{E}\) 在 \(\mathrm{L}_{\mathrm{K}}^{p}(\mathrm{X},\mu)\) 中稠密（INT, IV, p. 162, §4, 第 10 目, 推论 1），存在一个有限集 \(\mathrm{F}'\)（含于 \(\mathcal{E}\)），使得 F 的每个元素都以至多 \(\varepsilon\) 的距离接近 \(\mathrm{F}'\) 的某个元素。考虑由诸测度非零的集合所构成的有限分割 \(\pi\)，映射 \(x\mapsto (f(x))_{f\in \mathrm{F}'}\) 在每个这样的集合上取一个给定的值，便可看出存在一个元素 \(\pi\)（属于 \(\mathcal{P}\)），使得 \(u_{\pi}(f) = f\) 对一切 \(f\)（含于 \(\mathrm{F}'\)）成立。因此投影子 \(u_{\pi}\) 具有所要的性质。

命题 14. --- 设 X 是局部紧拓扑空间，\(\mu\) 是 X 上的一个正测度，\(p \in [1, +\infty]\)。空间 \(\mathrm{L}_{\mathrm{K}}^{p}(\mathrm{X}, \mu)\) 是逼近空间。

若 \(p = +\infty\)，则空间 \(\mathrm{L}_{\mathbf{C}}^{\infty}(\mathrm{X},\mu)\) 是交换星代数（I, p. 103, 例 4），因而是逼近空间（命题 13 的推论），而由 III, p. 15 的注记 5，\(\mathrm{L}_{\mathbf{R}}^{\infty}(\mathrm{X},\mu)\) 亦然。

假设 p 有限。设 \(\mathrm{A}\subset\mathcal{L}^{\mathrm{f}}(\mathrm{L}_{\mathrm{K}}^{p}(\mathrm{X},\mu))\) 是范数 \(\leqslant1\) 的有限秩投影子所成的集合。由引理 6，\(\mathrm{L}_{\mathrm{K}}^{p}(\mathrm{X},\mu)\) 的恒等映射对逐点收敛拓扑附着于 A，于是命题由 III, p. 14 的注记 3 得出。

